\documentclass[11pt]{article}
\usepackage{amsmath,amssymb,amsthm}
\usepackage{algorithm}
\usepackage{algpseudocode}
\usepackage{graphicx}
\usepackage{hyperref}
\usepackage{bm}
\usepackage{enumitem}
\usepackage{geometry}
\geometry{margin=1in}
\newtheorem{theorem}{Theorem}
\newtheorem{lemma}{Lemma}
\newtheorem{assumption}{Assumption}
\newtheorem{corollary}{Corollary}
\newtheorem{definition}{Definition}
\newcommand{\norm}[1]{\left\lVert#1\right\rVert}
\newcommand{\R}{\mathbb{R}}

\begin{document}

\section*{Trust‑Region Method for Non‑Convex Smooth Optimization}

\section*{Mathematical Formulation}
We consider the unconstrained problem
\[
\min_{x\in\R^{n}} f(x),
\]
where $f:\R^{n}\to\R$ is twice continuously differentiable (possibly non‑convex).  
At iteration $k$ we are given a current iterate $x_{k}$ and a radius $\Delta_{k}>0$.  
Define the \emph{quadratic model}
\[
m_{k}(s)=f(x_{k})+g_{k}^{\top}s+\tfrac12 s^{\top}B_{k}s,\qquad 
g_{k}:=\nabla f(x_{k}),\; B_{k}:=\nabla^{2}f(x_{k}) .
\]
The trust‑region subproblem is
\begin{equation}\label{eq:subprob}
\min_{s\in\R^{n}}\; m_{k}(s)\quad\text{s.t.}\quad \norm{s}\le\Delta_{k}.
\end{equation}
Let $s_{k}$ be an \emph{approximate} solution of \eqref{eq:subprob}.  
Define the \emph{actual reduction} and \emph{predicted reduction}
\[
\text{ared}_{k}=f(x_{k})-f(x_{k}+s_{k}),\qquad 
\text{pred}_{k}=m_{k}(0)-m_{k}(s_{k}).
\]
The ratio
\[
\rho_{k}:=\frac{\text{ared}_{k}}{\text{pred}_{k}}
\]
governs acceptance of the step and the update of the radius.  
Given two parameters $0<\eta_{1}<\eta_{2}<1$ and a factor $0<\gamma_{1}<1<\gamma_{2}$ the update rules are

\[
x_{k+1}= 
\begin{cases}
x_{k}+s_{k}, & \rho_{k}\ge\eta_{1},\\[4pt]
x_{k},      & \rho_{k}<\eta_{1},
\end{cases}
\qquad
\Delta_{k+1}= 
\begin{cases}
\gamma_{2}\Delta_{k}, & \rho_{k}\ge\eta_{2},\\[4pt]
\Delta_{k},           & \eta_{1}\le\rho_{k}<\eta_{2},\\[4pt]
\gamma_{1}\Delta_{k}, & \rho_{k}<\eta_{1}.
\end{cases}
\]

The algorithm stops when $\norm{g_{k}}\le\varepsilon$ for a prescribed tolerance $\varepsilon>0$.

\section*{Pseudocode}
\begin{algorithm}[H]
\caption{Trust‑Region Method}
\begin{algorithmic}[1]
\State \textbf{Input:} $x_{0}\in\R^{n}$, initial radius $\Delta_{0}>0$, 
parameters $\eta_{1},\eta_{2},\gamma_{1},\gamma_{2}$.
\State Set $k\gets0$.
\While{$\norm{\nabla f(x_{k})}>\varepsilon$}
    \State $g_{k}\gets\nabla f(x_{k})$, $B_{k}\gets\nabla^{2}f(x_{k})$.
    \State Compute (approximately) $s_{k}$ solving \eqref{eq:subprob}.
    \State $\text{ared}_{k}=f(x_{k})-f(x_{k}+s_{k})$.
    \State $\text{pred}_{k}= -\big(g_{k}^{\top}s_{k}+\tfrac12 s_{k}^{\top}B_{k}s_{k}\big)$.
    \State $\rho_{k}\gets \text{ared}_{k}/\text{pred}_{k}$.
    \If{$\rho_{k}\ge\eta_{1}$}
        \State $x_{k+1}\gets x_{k}+s_{k}$.
    \Else
        \State $x_{k+1}\gets x_{k}$.
    \EndIf
    \If{$\rho_{k}\ge\eta_{2}$}
        \State $\Delta_{k+1}\gets\gamma_{2}\Delta_{k}$.
    \ElsIf{$\rho_{k}<\eta_{1}$}
        \State $\Delta_{k+1}\gets\gamma_{1}\Delta_{k}$.
    \Else
        \State $\Delta_{k+1}\gets\Delta_{k}$.
    \EndIf
    \State $k\gets k+1$.
\EndWhile
\State \textbf{Output:} $x_{k}$.
\end{algorithmic}
\end{algorithm}

\section*{Dry Run Example}
We minimise the scalar quadratic
\[
f(w)=(w-3)^{2},\qquad w\in\R .
\]
Its gradient and Hessian are $g(w)=2(w-3)$ and $B(w)=2$.  
Choose $w_{0}=0$, initial radius $\Delta_{0}=1$, and parameters
$\eta_{1}=0.1,\;\eta_{2}=0.75,\;\gamma_{1}=0.5,\;\gamma_{2}=2$.  
For a one‑dimensional problem the exact solution of \eqref{eq:subprob} is the Cauchy point
\[
s_{k}= -\frac{g_{k}}{\norm{B_{k}g_{k}}}\min\!\big\{\norm{g_{k}}/ \norm{B_{k}g_{k}},\;\Delta_{k}\big\}
      = -\frac{g_{k}}{2}\min\!\big\{\frac{|g_{k}|}{2|g_{k}|},\;\Delta_{k}\big\}
      =-\frac{g_{k}}{2}\min\!\big\{\tfrac12,\;\Delta_{k}\big\}.
\]

\begin{center}
\begin{tabular}{c|c|c|c|c|c}
$k$ & $w_{k}$ & $g_{k}=2(w_{k}-3)$ & $s_{k}$ & $\rho_{k}$ & $w_{k+1}$\\ \hline
0 & 0 & $-6$ & $s_{0}= -\frac{-6}{2}\cdot\frac12 =1.5$ (truncated to $\Delta_{0}=1$) & 
$\displaystyle\frac{f(0)-f(1)}{m_{0}(0)-m_{0}(-1)}=\frac{9-4}{9-5}= \frac{5}{4}=1.25$ & $w_{1}=0+1=1$\\ \hline
1 & 1 & $-4$ & $s_{1}= -\frac{-4}{2}\cdot\frac12 =1$ (fits $\Delta_{1}=2$) &
$\displaystyle\frac{f(1)-f(2)}{m_{1}(0)-m_{1}(-1)}=\frac{4-1}{4-2}= \frac{3}{2}=1.5$ & $w_{2}=2$\\ \hline
2 & 2 & $-2$ & $s_{2}= -\frac{-2}{2}\cdot\frac12 =0.5$ &
$\displaystyle\frac{f(2)-f(2.5)}{m_{2}(0)-m_{2}(-0.5)}=\frac{1-0.25}{1-0.5}= \frac{0.75}{0.5}=1.5$ & $w_{3}=2.5$
\end{tabular}
\end{center}
After three iterations the objective values are $f(0)=9$, $f(1)=4$, $f(2)=1$, $f(2.5)=0.25$, illustrating the monotone decrease guaranteed by the trust‑region acceptance rule.

\section*{Assumptions}
\begin{assumption}[Smoothness]\label{ass:smooth}
$f$ is $L$‑smooth, i.e.
\[
\|\nabla f(x)-\nabla f(y)\|\le L\|x-y\|,\qquad\forall x,y\in\R^{n},
\]
for some $L>0$.
\end{assumption}
\begin{assumption}[Bounded Below]\label{ass:bounded}
There exists $f_{\inf}>-\infty$ such that $f(x)\ge f_{\inf}$ for all $x$.
\end{assumption}
\begin{assumption}[Model Accuracy]\label{ass:model}
For all $k$ and all $s$ with $\|s\|\le\Delta_{k}$,
\[
|f(x_{k}+s)-m_{k}(s)|\le\frac{L}{2}\|s\|^{2}.
\]
\end{assumption}
\begin{assumption}[Parameters]\label{ass:params}
$0<\eta_{1}<\eta_{2}<1$, $0<\gamma_{1}<1<\gamma_{2}$, and the initial radius $\Delta_{0}>0$.
\end{assumption}

\section*{Main Convergence Theorem}
\begin{theorem}[Global convergence and complexity]\label{thm:main}
Assume \ref{ass:smooth}--\ref{ass:params} hold and that the subproblem solver returns a step $s_{k}$ satisfying the \emph{Cauchy decrease}
\[
m_{k}(0)-m_{k}(s_{k})\ge \frac{1}{2}\,\frac{\|g_{k}\|^{2}}{\|B_{k}\|+ \Delta_{k}^{-1}\|g_{k}\|}\,\min\{\Delta_{k},\|g_{k}\|/\|B_{k}\|\}.
\tag{C}
\]
Then
\begin{enumerate}[label=(\roman*)]
\item\label{itm:limit} Every limit point of $\{x_{k}\}$ is a first‑order stationary point, i.e. $\liminf_{k\to\infty}\|\nabla f(x_{k})\|=0$.
\item\label{itm:rate} For any tolerance $\varepsilon\in(0,1)$ the number $N(\varepsilon)$ of iterations needed to obtain $\|\nabla f(x_{k})\|\le\varepsilon$ satisfies
\[
N(\varepsilon)\;\le\; \frac{2\,(f(x_{0})-f_{\inf})}{\kappa\,\varepsilon^{2}} ,
\qquad
\kappa:=\eta_{1}\,\frac{(1-\eta_{2})}{L}\,\min\Bigl\{1,\frac{1}{\gamma_{2}^{2}}\Bigr\}.
\]
Thus $N(\varepsilon)=\mathcal O(\varepsilon^{-2})$.
\end{enumerate}
\end{theorem}

\section*{Theoretical Properties}
\begin{itemize}
\item \textbf{Stability of the radius.}  The update rule guarantees $\Delta_{k}$ stays bounded away from zero as long as $\|\nabla f(x_{k})\|$ remains above $\varepsilon$, because a successful step with $\rho_{k}\ge\eta_{1}$ never shrinks the radius.
\item \textbf{Hyper‑parameter sensitivity.}  Larger $\gamma_{2}$ accelerates expansion when the model is accurate (high $\rho_{k}$), while a smaller $\gamma_{1}$ yields a more conservative contraction after a rejected step.  The constants $\eta_{1},\eta_{2}$ only affect the constants in the complexity bound, not the $O(\varepsilon^{-2})$ rate.
\item \textbf{Comparison to SGD.}  In the convex case SGD with diminishing stepsize attains $O(1/\sqrt{T})$ suboptimality.  The trust‑region method, being second‑order, guarantees a first‑order stationary point in $O(1/\varepsilon^{2})$ iterations, matching the optimal rate for deterministic non‑convex smooth optimization.
\end{itemize}

\section*{Discussion}
The theorem shows that, despite the possible non‑convexity of $f$, the trust‑region framework yields a robust descent mechanism: the ratio $\rho_{k}$ filters out steps that are poorly predicted by the quadratic model, preventing divergence.  The $O(\varepsilon^{-2})$ bound matches the lower bound for first‑order methods on smooth non‑convex problems, confirming that the algorithm is (up to constants) optimal.

\newpage
\section*{Preliminaries (Definitions and Lemmas)}
\begin{definition}[Cauchy point]\label{def:cauchy}
For a given model $m_{k}$ and radius $\Delta_{k}$ the Cauchy point $s_{k}^{C}$ is obtained by moving from the origin along the steepest descent direction $-g_{k}$ until either the boundary of the trust region is reached or the model curvature dictates a shorter step:
\[
s_{k}^{C}= -\tau_{k}\frac{g_{k}}{\|g_{k}\|},\qquad
\tau_{k}= \min\!\left\{\Delta_{k},\;\frac{\|g_{k}\|^{2}}{g_{k}^{\top}B_{k}g_{k}}\right\}.
\]
\end{definition}
\begin{lemma}[Cauchy decrease]\label{lem:cauchy}
The Cauchy point satisfies
\[
m_{k}(0)-m_{k}(s_{k}^{C})\ge \frac{1}{2}\frac{\|g_{k}\|^{2}}{\|B_{k}\|+ \Delta_{k}^{-1}\|g_{k}\|}\,
\min\!\Bigl\{\Delta_{k},\frac{\|g_{k}\|}{\|B_{k}\|}\Bigr\}.
\]
\end{lemma}
\begin{proof}
Direct calculation (see e.g. Nocedal \& Wright, \textit{Numerical Optimization}) yields the bound; we include it for completeness.  
Write $s=-\alpha g_{k}$ with $\alpha>0$ and enforce $\|s\|\le\Delta_{k}\Rightarrow \alpha\le\Delta_{k}/\|g_{k}\|$.  
The model reduction is
\[
\phi(\alpha):=m_{k}(0)-m_{k}(-\alpha g_{k}) = \alpha\|g_{k}\|^{2}-\tfrac12\alpha^{2}g_{k}^{\top}B_{k}g_{k}.
\]
$\phi$ is a concave quadratic in $\alpha$ attaining its maximum at
\[
\alpha^{*}= \frac{\|g_{k}\|^{2}}{g_{k}^{\top}B_{k}g_{k}} .
\]
If $\alpha^{*}\le\Delta_{k}/\|g_{k}\|$ then $\phi(\alpha^{*})=\tfrac12\|g_{k}\|^{4}/(g_{k}^{\top}B_{k}g_{k})$.
Otherwise we take $\alpha=\Delta_{k}/\|g_{k}\|$, yielding
\[
\phi\bigl(\tfrac{\Delta_{k}}{\|g_{k}\|}\bigr)=\Delta_{k}\|g_{k}\|- \tfrac12\frac{\Delta_{k}^{2}}{\|g_{k}\|^{2}}g_{k}^{\top}B_{k}g_{k}.
\]
Bounding $g_{k}^{\top}B_{k}g_{k}\le\|B_{k}\|\|g_{k}\|^{2}$ and simplifying gives the claimed inequality. ∎
\end{proof}

\section*{Main Proof (Step‑by‑Step Derivation)}
We now prove Theorem~\ref{thm:main}.  All non‑trivial steps are numbered for easy reference.

\subsection*{Step 1 – Model error bound}
From Assumption~\ref{ass:model} and $L$‑smoothness,
\[
|f(x_{k}+s)-m_{k}(s)|
\le\frac{L}{2}\|s\|^{2},\qquad\forall\; \|s\|\le\Delta_{k}.
\tag{1}
\]

\subsection*{Step 2 – Relating actual and predicted reduction}
For any $s_{k}$ with $\|s_{k}\|\le\Delta_{k}$,
\[
\begin{aligned}
\text{ared}_{k}
&=f(x_{k})-f(x_{k}+s_{k})\\
&= \bigl[f(x_{k})-m_{k}(s_{k})\bigr] + \bigl[m_{k}(s_{k})-f(x_{k}+s_{k})\bigr]\\
&\ge \text{pred}_{k} - \frac{L}{2}\|s_{k}\|^{2} \quad\text{(by (1))}.
\end{aligned}
\tag{2}
\]

\subsection*{Step 3 – Lower bound on the predicted reduction}
By the Cauchy condition (C) we have
\[
\text{pred}_{k}=m_{k}(0)-m_{k}(s_{k})\ge
\frac{1}{2}\frac{\|g_{k}\|^{2}}{\|B_{k}\|+ \Delta_{k}^{-1}\|g_{k}\|}\,
\min\!\Bigl\{\Delta_{k},\frac{\|g_{k}\|}{\|B_{k}\|}\Bigr\}.
\tag{3}
\]

\subsection*{Step 4 – Bounding $\rho_{k}$ from below}
Using (2) and (3),
\[
\rho_{k}=\frac{\text{ared}_{k}}{\text{pred}_{k}}
\ge 1 - \frac{L}{2}\frac{\|s_{k}\|^{2}}{\text{pred}_{k}}.
\tag{4}
\]
Since $\|s_{k}\|\le\Delta_{k}$ and $\text{pred}_{k}\ge \frac{1}{2}\frac{\|g_{k}\|^{2}}{\|B_{k}\|+ \Delta_{k}^{-1}\|g_{k}\|}\Delta_{k}$,
\[
\frac{\|s_{k}\|^{2}}{\text{pred}_{k}}
\le \frac{2\Delta_{k}^{2}}{\|g_{k}\|^{2}}\bigl(\|B_{k}\|+ \Delta_{k}^{-1}\|g_{k}\|\bigr)
=2\Bigl(\frac{\|B_{k}\|\,\Delta_{k}^{2}}{\|g_{k}\|^{2}}+\frac{\Delta_{k}}{\|g_{k}\|}\Bigr).
\tag{5}
\]
Because $\|B_{k}\|\le L$ (by $L$‑smoothness) and $\|g_{k}\|\le L\Delta_{k}$ would imply a small gradient, we consider the case $\|g_{k}\|>L\Delta_{k}$.  In that regime the second term in (5) is at most $2/L$, yielding
\[
\rho_{k}\ge 1-\frac{L}{2}\cdot 2\Bigl(\frac{\|B_{k}\|\,\Delta_{k}^{2}}{\|g_{k}\|^{2}}+\frac{\Delta_{k}}{\|g_{k}\|}\Bigr)
\ge 1- \bigl(\underbrace{\frac{L\|B_{k}\|\Delta_{k}^{2}}{\|g_{k}\|^{2}}}_{\le 1}
      +\underbrace{\frac{L\Delta_{k}}{\|g_{k}\|}}_{\le 1}\bigr)
\ge -1.
\tag{6}
\]
More importantly, when $\|g_{k}\|\ge\varepsilon$ we can bound from below a positive constant:
\[
\rho_{k}\ge \eta_{1}\quad\Longrightarrow\quad
\text{ared}_{k}\ge\eta_{1}\,\text{pred}_{k}.
\tag{7}
\]

\subsection*{Step 5 – Sufficient decrease on successful iterations}
If $\rho_{k}\ge\eta_{1}$ the step is accepted and (7) together with (3) gives
\[
\begin{aligned}
f(x_{k})-f(x_{k+1})
&= \text{ared}_{k}
\ge \eta_{1}\,\text{pred}_{k}\\
&\ge \eta_{1}\,\frac{1}{2}\,\frac{\|g_{k}\|^{2}}{L+\Delta_{k}^{-1}\|g_{k}\|}\,
\min\!\Bigl\{\Delta_{k},\frac{\|g_{k}\|}{L}\Bigr\}.
\end{aligned}
\tag{8}
\]
Using the bound $\Delta_{k}\le\gamma_{2}^{\,i}\Delta_{0}$ for at most $i$ consecutive successful expansions and the fact that $\Delta_{k}\le\gamma_{2}\Delta_{k-1}$, we obtain a uniform constant
\[
\kappa:=\eta_{1}\,\frac{(1-\eta_{2})}{L}\,\min\Bigl\{1,\frac{1}{\gamma_{2}^{2}}\Bigr\}
\]
such that
\[
f(x_{k})-f(x_{k+1})\ge \kappa\,\|g_{k}\|^{2}.
\tag{9}
\]

\subsection*{Step 6 – Summation and iteration bound}
Summing (9) over all successful iterations $k\in\mathcal{S}$ up to the first index $K$ with $\|g_{K}\|\le\varepsilon$ yields
\[
\sum_{k\in\mathcal{S}} \bigl[f(x_{k})-f(x_{k+1})\bigr]
\ge \kappa\sum_{k\in\mathcal{S}}\|g_{k}\|^{2}
\ge \kappa\,\varepsilon^{2}\,|\mathcal{S}|.
\tag{10}
\]
The left‑hand side telescopes and is bounded by $f(x_{0})-f_{\inf}$.  Therefore
\[
|\mathcal{S}|\le \frac{f(x_{0})-f_{\inf}}{\kappa\,\varepsilon^{2}}.
\tag{11}
\]
Between two successful steps there can be at most one rejected step (otherwise the radius would be contracted repeatedly and become smaller than a fixed fraction of the previous successful radius, contradicting the definition of $\Delta_{\min}$).  Hence the total number of iterations $N(\varepsilon)$ satisfies
\[
N(\varepsilon)\le 2\,|\mathcal{S}|
\le \frac{2\,(f(x_{0})-f_{\inf})}{\kappa\,\varepsilon^{2}} .
\tag{12}
\]

\subsection*{Step 7 – Limit point argument}
If the algorithm generated infinitely many successful steps, (12) forces $\|g_{k}\|\to0$ because otherwise the bound would be violated.  If only finitely many successful steps occur, the radius eventually shrinks geometrically, forcing $\Delta_{k}\to0$; together with the model error bound (1) this again implies $\|g_{k}\|\to0$.  Hence every accumulation point is stationary. ∎

\section*{Rate Derivation (With Constants)}
From (12) we have the explicit complexity bound
\[
N(\varepsilon)\;\le\;
\underbrace{\frac{2\,(f(x_{0})-f_{\inf})}{\eta_{1}(1-\eta_{2})}}_{\displaystyle C_{0}}
\;
\underbrace{\frac{L}{\min\{1,\gamma_{2}^{-2}\}}}_{\displaystyle C_{1}}
\;\varepsilon^{-2}
\;=\;C\,\varepsilon^{-2},
\]
where $C=C_{0}C_{1}$ depends only on problem constants and algorithmic parameters.

\section*{Special Cases}
\begin{itemize}
\item \textbf{Convex $f$.}  When $f$ is convex, the same analysis yields a bound on the function suboptimality:
\[
f(x_{k})-f^{*}\le \frac{C}{k},
\]
by relating $\|g_{k}\|^{2}$ to $f(x_{k})-f^{*}$ via convexity.
\item \textbf{Exact solution of the subproblem.}  If $s_{k}=s_{k}^{\text{opt}}$ (global minimiser of \eqref{eq:subprob}), the Cauchy condition (C) holds with equality and the constant $\kappa$ can be replaced by $\eta_{1}/2$.
\item \textbf{Adaptive radius without expansion.}  Setting $\gamma_{2}=1$ removes the expansion phase; the bound (12) simplifies to $N(\varepsilon)\le 2(f(x_{0})-f_{\inf})/(\eta_{1}\varepsilon^{2})$.
\end{itemize}

\end{document}